% ---------------------------------------------------------------------
%  NỘI DUNG CHƯƠNG 2 — CÀI ĐẶT VÀ THỬ NGHIỆM
% ---------------------------------------------------------------------
\section{Cài đặt và thử nghiệm}
\label{sec:caidat}

\begin{muctieu}
Phần này mô tả công nghệ, các thành phần đã triển khai và cách kiểm thử
những luồng nghiệp vụ chính của hệ thống.
\end{muctieu}

\subsection{Môi trường triển khai}

Dự án được tổ chức thành frontend web, backend API, cơ sở dữ liệu và cấu
hình Docker Compose. Frontend dùng React 18 và TypeScript 5; backend dùng
Node.js với Express 5; PostgreSQL 16 là dịch vụ cơ sở dữ liệu trong cấu
hình container. Các phiên bản thư viện chính được liệt kê ở
Bảng~\ref{tab:congcu}. Dockerfile của frontend dùng bước build Vite và Nginx
để phục vụ tệp tĩnh; backend chạy trên ảnh Node.js Alpine.

Mã nguồn được kiểm tra tại thư mục cục bộ
\code{E:\textbackslash AI-EXAMHUB-DAU}. Cấu hình môi trường được cung cấp qua
các tệp \code{.env.example}; khóa dịch vụ và mật khẩu dùng khi triển khai
thực tế cần được thay bằng giá trị riêng, không đưa vào báo cáo hoặc kho mã
nguồn. Hướng dẫn chạy từ mã nguồn được trình bày ở Phụ lục.

\subsubsection{Cấu trúc mã nguồn}
Các thành phần được tách theo trách nhiệm. \code{frontend/web/src/pages/}
chứa các màn hình nghiệp vụ; \code{frontend/web/src/services/api.ts} là lớp
gọi backend. \code{backend/api/src/app.js} khởi tạo ứng dụng Express và các
route; \code{backend/api/src/routes/} chứa các nhóm API. SQL nền tảng nằm
trong \code{database/schema.sql}; các thay đổi theo phiên bản được đặt tại
\code{database/migrations/}.

\begin{table}[H]
\centering
\caption{Phân chia trách nhiệm theo thư mục mã nguồn.}
\label{tab:modules}
\begin{tabularx}{\textwidth}{>{\bfseries\color{daunavy}}p{0.29\textwidth} X}
\toprule
\rowcolor{gray100}Đường dẫn & Trách nhiệm trong hệ thống \\
\midrule
\code{frontend/web/src/pages/} & Màn hình tạo đề, câu hỏi, thi, lớp học, người dùng, chấm và thống kê. \\
\code{frontend/web/src/services/api.ts} & Tạo yêu cầu HTTP, giữ kiểu dữ liệu API và quản lý thông tin xác thực ở phía client. \\
\code{backend/api/src/app.js} & Khởi tạo Express, tích hợp phân tích tài liệu và đăng ký route. \\
\code{backend/api/src/routes/} & Xác thực, câu hỏi, đề thi, phòng thi, lớp học, người dùng, bài nộp và phân tích. \\
\code{backend/api/src/middlewares/} & Xác thực JWT và kiểm tra vai trò. \\
\code{database/schema.sql} & Định nghĩa bảng, khoá, ràng buộc và chỉ mục chính. \\
\code{database/migrations/} & Chuyển đổi schema cho các phiên bản chức năng. \\
\code{docker-compose.yml} & Khai báo dịch vụ phát triển và mạng/kết nối giữa các container. \\
\bottomrule
\end{tabularx}
\end{table}

\subsubsection{Biến cấu hình và phụ thuộc ngoài}
Bảng~\ref{tab:config} mô tả nhóm cấu hình theo mục đích, không nêu giá trị
bí mật. Ứng dụng có thể chạy phần sinh dự phòng mà không cấu hình API mô
hình; chức năng Google OAuth cũng phụ thuộc thông tin client do nhà triển
khai cấp. Hướng dẫn Node.js và Express được dùng để tham chiếu hành vi môi
trường chạy và định tuyến \cite{NodeDocs,ExpressDocs}.
\begin{table}[H]
\centering
\caption{Nhóm cấu hình cần thiết khi chạy hệ thống.}
\label{tab:config}
\begin{tabularx}{\textwidth}{>{\bfseries\color{daunavy}}p{0.22\textwidth}p{0.33\textwidth}X}
\toprule
\rowcolor{gray100}Nhóm biến & Ví dụ tên biến & Mục đích / lưu ý \\
\midrule
API và xác thực & \code{PORT}, \code{JWT\_SECRET}, \code{JWT\_EXPIRES\_IN} & Cổng API, ký và thời hạn token; secret phải duy nhất, không dùng giá trị dev khi triển khai. \\
Cơ sở dữ liệu & \code{DB\_HOST}, \code{DB\_PORT}, \code{DB\_NAME}, \code{DB\_USER}, \code{DB\_PASSWORD} & Kết nối PostgreSQL; mật khẩu lấy từ môi trường an toàn. \\
Giao diện & \code{VITE\_API\_URL} & Địa chỉ gốc của API được đóng vào frontend. \\
LLM tùy chọn & \code{OPENAI\_API\_KEY}, \code{OPENAI\_MODEL} & Gọi API sinh nội dung; khóa không được chuyển xuống trình duyệt. \\
Đăng nhập Google & \code{GOOGLE\_CLIENT\_ID}, \code{GOOGLE\_CLIENT\_SECRET}, \code{GOOGLE\_REDIRECT\_URI} & Chỉ hoạt động khi được đăng ký/cấu hình đúng tại nhà cung cấp. \\
Nguồn gọi API & \code{CORS\_ORIGIN}, \code{FRONTEND\_URL} & Chính sách CORS và URL chuyển hướng OAuth. \\
\bottomrule
\end{tabularx}
\end{table}

\subsubsection{Sơ đồ triển khai các thành phần}
Hình~\ref{fig:deployment} khái quát luồng kết nối của phiên bản hiện tại.
Trình duyệt tải giao diện tĩnh do Nginx phục vụ, sau đó gửi yêu cầu đến API
Express; backend truy cập PostgreSQL và chỉ gọi nhà cung cấp LLM bên ngoài
khi đã cấu hình khóa. Các thành phần frontend, backend và cơ sở dữ liệu được
khai báo trong Docker Compose \cite{DockerComposeDocs}; mô hình không hàm ý
dịch vụ AI riêng biệt chạy trong Compose.
\begin{figure}[H]
\centering
\resizebox{0.95\textwidth}{!}{%
\begin{tikzpicture}[node distance=10mm and 13mm]
  \node[objbox] (browser) {Trình duyệt\\React};
  \node[clsbox, right=of browser] (web) {Nginx\\tệp giao diện};
  \node[clsbox, right=of web] (api) {Express API\\nghiệp vụ};
  \node[clsbox, right=of api] (db) {PostgreSQL\\dữ liệu};
  \node[clsbox, below=of api] (llm) {Dịch vụ LLM\\tuỳ cấu hình};
  \draw[flowarr] (browser) -- node[above,font=\small] {HTTPS} (web);
  \draw[flowarr] (web) -- node[above,font=\small] {HTTP API} (api);
  \draw[flowarr] (api) -- node[above,font=\small] {SQL} (db);
  \draw[flowarr,dashed] (api) -- node[right,font=\small] {HTTPS} (llm);
\end{tikzpicture}}
\caption{Sơ đồ triển khai logic và các kết nối chính.}
\label{fig:deployment}
\end{figure}

\subsection{Các thành phần chính}

\subsubsection{Giao diện web}
Ứng dụng React có màn hình đăng nhập, trang tổng quan, ngân hàng câu hỏi,
trình tạo đề, duyệt đề, quản lý người dùng, chấm bài, phân tích, thi trực
tuyến và lớp học. Giao diện gọi các hàm tập trung trong dịch vụ API ở
\code{frontend/web/src/services/api.ts}. Một số trang điều phối luồng cụ thể
gồm \code{ExamBuilderPage.tsx}, \code{OnlineExamPage.tsx},
\code{QuestionBankPage.tsx} và \code{GradingPage.tsx}.

\subsubsection{API và nghiệp vụ}
Backend Express tổ chức endpoint theo nhóm xác thực, người dùng, câu hỏi,
đề thi, bài nộp, phân tích, phòng thi và lớp học. Middleware xác thực token
và kiểm tra vai trò được dùng để giới hạn thao tác. Endpoint phân tích tài
liệu tiếp nhận nội dung, gọi dịch vụ ngôn ngữ nếu cấu hình, kiểm tra đầu ra
và trả các câu hỏi đề xuất kèm thông tin bằng chứng. Nhóm endpoint thi trực
tuyến hỗ trợ tạo phòng, bắt đầu/tiếp tục lượt thi, lưu câu trả lời, nộp bài
và giám sát tiến độ.

\subsubsection{Lưu trữ và xử lý tài liệu}
Lược đồ PostgreSQL lưu người dùng, câu hỏi, đề thi, phòng/lượt thi, bài nộp,
kết quả và thành viên lớp. Các thay đổi lược đồ được quản lý bằng các tệp
migration trong \code{database/migrations/}. Tệp tải lên được nhận trong bộ
nhớ của backend để trích xuất nội dung; PDF dùng \code{pdf-parse}, DOC/DOCX
được chuyển qua Mammoth và TXT được giải mã UTF-8. Backend có nhánh OCR ảnh
dùng Tesseract.js với dữ liệu tiếng Anh, nhưng giao diện hiện không cung cấp
đầu vào ảnh. Báo cáo không giả định tài liệu đã được lưu trữ lâu dài hay
mọi tài liệu được trích xuất chính xác.

\subsubsection{Quản lý vòng đời đề thi}
Đề lưu trạng thái nghiệp vụ nháp, chờ duyệt, đã duyệt và phát hành cùng
metadata như thời lượng, tổng điểm, năm học, học kỳ và ma trận câu hỏi. Khi
lưu đề, backend kiểm tra số lượng, loại câu, điểm và thời lượng; các thao
tác cần nhiều câu lệnh sử dụng transaction để tránh lưu dở dữ liệu giữa
bảng đề và bảng liên kết. Câu hỏi thuộc đề được truy vấn qua
\code{exam\_questions} theo vị trí, thay vì dựa duy nhất vào một danh sách
ID không có ràng buộc quan hệ.

\subsubsection{Quản lý phiên làm bài}
Khi sinh viên vào phòng, API tạo hoặc lấy lại lượt thi. API lưu đáp án theo
ID câu hỏi, danh sách câu đánh dấu, chỉ số câu hiện tại, thời hạn và thời
điểm nộp. Khi nộp, bài thi trắc nghiệm được tự chấm; đề tự luận có thể yêu
cầu chấm riêng. Trước khi trả câu hỏi cho sinh viên, route loại các trường
đáp án đúng, đáp án tham khảo và bằng chứng nội bộ. Cách tách thông tin
đáp án là kiểm soát quan trọng vì ẩn đáp án chỉ ở giao diện không đủ để
bảo vệ bài thi.

\subsubsection{Thống kê và chấm điểm}
Trang thống kê hiển thị số liệu theo quyền người dùng; quản trị viên có
phạm vi tổng thể, giảng viên tập trung dữ liệu liên quan đề/lớp của mình,
còn sinh viên xem lượt dự thi. Chấm thủ công nhận điểm và phản hồi, sau đó
lưu vào bảng \code{results}; điểm được kiểm tra không âm và không vượt tổng
điểm của đề. Giao diện có nội dung mô tả “AI feedback”, nhưng backend được
kiểm tra trong đồ án không chứng minh đã có pipeline tự chấm tự luận bằng
mô hình; báo cáo chỉ ghi nhận luồng chấm do giảng viên thực hiện.

\subsection{Giao diện và các luồng chính}

Các màn hình triển khai bao phủ hai nhóm luồng. Giảng viên có thể quản lý
câu hỏi, phân tích tài liệu, chọn câu để tạo đề, gửi duyệt, phát hành và
tạo phòng thi. Sinh viên đăng nhập, tham gia bằng mã, làm bài, lưu tiến độ,
tiếp tục lượt thi và xem kết quả. Trang giám sát của giảng viên tải lại
tiến độ theo chu kỳ; trang lớp học quản lý mã tham gia và danh sách thành
viên. Báo cáo hiện không kèm ảnh chụp giao diện vì thư mục dự án không có
ảnh chụp màn hình sản phẩm phù hợp để làm minh chứng.

\subsubsection{Minh hoạ bố cục khu vực tạo đề}
Hình~\ref{fig:wireframe} là sơ đồ bố cục khái niệm được dựng để giải thích
các vùng chức năng, không phải ảnh chụp từ trình duyệt. Nó chỉ ra các bước
người dùng cần quan sát: cấu hình đề, nạp nguồn, chọn câu hỏi và xác nhận
trước khi lưu.
\begin{figure}[H]
\centering
\resizebox{0.92\textwidth}{!}{%
\begin{tikzpicture}[x=1cm,y=1cm]
  \draw[rounded corners=3pt, line width=1pt, color=daunavy] (0,0) rectangle (15,8);
  \fill[lightnavy] (0,7) rectangle (15,8);
  \node[anchor=west,color=daunavy,font=\bfseries] at (0.4,7.5) {AI EXAMHUB DAU \hspace{1em} Tạo đề thi};
  \fill[gray100] (0.4,0.5) rectangle (4.4,6.6);
  \node[anchor=west,font=\bfseries,color=daunavy] at (0.7,6.15) {1. Cấu hình đề};
  \node[anchor=west,align=left] at (2.4,5.25) {Tên đề / môn học\\Thời lượng / tổng điểm\\Loại câu / độ khó\\Mức độ nhận thức};
  \draw[color=gray300] (0.7,3.8) -- (4.1,3.8);
  \node[anchor=west,align=left] at (2.4,3.0) {Chọn tài liệu nguồn\\Tải PDF / DOCX / TXT};
  \node[anchor=west,font=\bfseries,color=daunavy] at (0.7,1.35) {Phân tích tài liệu};
  \draw[rounded corners=2pt,color=daumaroon] (4.8,4.4) rectangle (14.5,6.6);
  \node[anchor=west,font=\bfseries,color=daunavy] at (5.2,6.15) {2. Nội dung nguồn và trích dẫn};
  \node[anchor=west,align=left] at (9.6,5.35) {Xem văn bản đã trích xuất\\Đối chiếu đoạn bằng chứng\\Điều chỉnh thông tin môn học};
  \draw[rounded corners=2pt,color=daumaroon] (4.8,1.0) rectangle (14.5,4.0);
  \node[anchor=west,font=\bfseries,color=daunavy] at (5.2,3.55) {3. Danh sách câu hỏi đề xuất};
  \node[anchor=west,align=left] at (9.6,2.4) {Câu hỏi trắc nghiệm / tự luận\\Mức độ, đáp án và căn cứ\\Chọn, bỏ hoặc hiệu chỉnh câu};
  \fill[lightwarm] (11.4,0.3) rectangle (14.5,0.8);
  \node[anchor=center,font=\bfseries,color=daumaroon] at (12.95,0.55) {Lưu đề};
\end{tikzpicture}}
\caption{Wireframe khái niệm cho các khu vực của quy trình tạo đề.}
\label{fig:wireframe}
\end{figure}

\subsubsection{Trạng thái thông báo và lỗi}
Trong quy trình tạo tài liệu, trạng thái giao diện cần phân biệt đang tải,
trích xuất xong, không có nội dung đọc được, không sinh được ứng viên, đã
sinh được ứng viên và lỗi dịch vụ. Người dùng phải biết liệu câu hỏi thuộc
nhánh LLM hay nhánh trích xuất dự phòng nếu hệ thống trả trường đó. Lỗi
không nên hiển thị dưới dạng thông báo thành công, và thao tác thử lại
không được tạo bản ghi trùng ngoài ý muốn. Những tiêu chí này được đưa vào
ma trận kiểm thử đề xuất ở Phụ lục.

\subsection{Kiểm thử}

Backend có bộ kiểm thử tích hợp tại \code{backend/api/tests/api.test.js}.
Trong lần chạy phục vụ báo cáo, toàn bộ 35 bài kiểm thử đạt, không có bài
thất bại. Các trường hợp tập trung vào đăng nhập và phân quyền, CRUD tài
khoản/câu hỏi/đề thi, phân tích tài liệu, sinh câu hỏi khi không cấu hình
mô hình ngoài, phòng thi trực tuyến, tiếp tục bài làm, tự chấm trắc nghiệm
và quản lý lớp học.

Frontend được kiểm tra bằng lệnh \code{npm run build}; TypeScript hoàn tất
kiểm tra kiểu và Vite tạo gói production thành công. Kết quả này xác nhận
build và các hành vi được kiểm thử, không phải đánh giá độ chính xác của AI,
độ chính xác OCR, tải đồng thời hay thời gian phản hồi.

\subsubsection{Tổ chức kiểm thử}
Bộ kiểm thử API dùng Node.js test runner và chạy ứng dụng trên cổng động
trong cùng tiến trình. Các test tạo tài khoản có vai trò, gửi HTTP request
đến route và truy vấn PostgreSQL để xác minh dữ liệu. Phần lớn kiểm thử
được tích hợp với database đã cấu hình trong môi trường này; do đó, người
chạy lại cần database khả dụng và phải bảo đảm test không chạy nhầm vào
dữ liệu production. Tài liệu của Node.js mô tả API test runner và cách
khai báo test bất đồng bộ \cite{NodeTestDocs}.

\subsubsection{Kiểm soát an toàn dữ liệu}
Hệ thống chứa thông tin định danh tài khoản, câu trả lời và lịch sử thi.
Nguyên tắc cần áp dụng gồm chỉ truy xuất dữ liệu theo quyền sở hữu, tránh
ghi token/khóa vào log, mã hoá đường truyền khi triển khai, sao lưu có kiểm
soát và quy định thời hạn lưu dữ liệu. Khi gửi tài liệu cho nhà cung cấp
LLM, cần thông báo rõ đối tượng dữ liệu, mục đích và cơ chế lưu/huỷ. Các
nguyên tắc này phù hợp với khung quản trị AI và an ninh ứng dụng
\cite{NISTAIRMF,OWASPTop10,OWASPAPI,OWASPLLM}; chúng chưa được xem là đã xác
minh nếu chưa có kiểm toán chuyên biệt.
